% Source: source/普通高考/2012/2012安徽文.pdf, page 1, question 6.
% Program-flow topology from the original PDF:
%   start -> init -> test; test --是--> update -> return to the main stem;
%   test --否--> output -> finish.
%   The return arrow is independent and points left into the stem above test;
%   the stem's only downward arrow enters test.

\begin{tikzpicture}[
  >=Stealth,
  x=1cm,
  y=0.82cm,
  line width=0.45pt,
  line join=round,
  line cap=round,
  font=\normalsize,
  flowtext/.style={anchor=center, align=center,
    execute at begin node={\setlength{\parindent}{0pt}}},
  every node/.style={flowtext},
  flowbox/.style={flowtext, draw, minimum height=0.68cm,
    inner xsep=4pt, inner ysep=2.5pt},
  branchlabel/.style={flowtext, inner sep=0pt},
  terminal/.style={flowbox, rounded corners=0.18cm, minimum width=1.35cm},
  process/.style={flowbox},
  decision/.style={flowbox, diamond, aspect=3.0, minimum width=3.55cm,
    minimum height=0.95cm},
  io/.style={flowbox, trapezium, trapezium left angle=72,
    trapezium right angle=108, minimum width=2.15cm, minimum height=0.76cm},
]
  \path[use as bounding box] (-1.95,2.90) rectangle (5.75,9.65);

  \node[terminal] (start) at (0,9.30) {开始};
  \node[process, minimum width=2.90cm] (init) at (0,7.92)
    {$x=1,y=1$};
  \node[decision] (test) at (0,6.30) {$x\leqslant 4?$};
  \node[process, minimum width=3.65cm, minimum height=0.78cm]
    (update) at (4.05,6.30) {$x=2x,y=y+1$};
  \node[io] (output) at (0,4.72) {输出 $y$};
  \node[terminal] (finish) at (0,3.35) {结束};

  % Main sequence; the stem segment above test has no arrow, and the
  % single downward arrow is placed immediately before the decision node.
  \draw[->] (start.south) -- (init.north);
  \coordinate (loopJoin) at (0,7.15);
  \draw (init.south) -- (loopJoin);
  \draw[->] (loopJoin) -- (test.north);
  \draw[->] (test.south) -- node[branchlabel, right=2pt] {否} (output.north);
  \draw[->] (output.south) -- (finish.north);

  % Affirmative branch and its independent return arrow.
  \draw[->] (test.east) -- node[branchlabel, above=2pt] {是} (update.west);
  \draw[->] (update.north) -- (4.05,7.15) -- (loopJoin);
\end{tikzpicture}
